\section{Phase-E Extensions Toward NeurIPS Main Track}
\label{sec:phaseE}

Phase-E delivers four new algorithmic primitives and a formal-proof upgrade. The pattern across these results is consistent with the synthesis of §\ref{sec:phaseD}: PRICE-RL is a strong, general-purpose primitive on short combinatorial / deceptive / small-budget / diverse-set / reward-hacking benchmarks, and it integrates cleanly with locality / multi-objective / token-level extensions while preserving Theorem~\ref{thm:exact} exactly. Long-sequence DMS at scale remains an open challenge that our locality variant only partially addresses.

\subsection{Locality-aware PRICE-RL}
\label{sec:e18}

We augment PRICE-RL with a per-candidate Hamming-distance kernel $w_{\mathrm{loc}}(x; \mathcal C, r) = \exp(-d_H(x, \mathcal C)/r)$, where $\mathcal C$ is the top-$k$ best-fitness centres in the labelled history. The locality-weighted gradient preserves Theorem~\ref{thm:exact} (see Appendix~\ref{app:locality}).

\begin{table}[t]
\centering
\caption{E18 --- Locality-aware PRICE-RL on long-sequence DMS, 4{,}000-query budget. Locality narrows the gap on Enhancer (within 1.3$\times$ of AdaLead) but does not close it on TEM-1. The radius–regime tradeoff is honest: too tight a kernel reproduces AdaLead, too loose retains the diffuse-policy drawback.}
\label{tab:e18}
\begin{tabular}{lcc}
\toprule
Dataset & AdaLead & PRICE-RL (locality) \\
\midrule
TEM-1 (L=287, $r=11.5$)        & $\mathbf{65.76 \pm 4.6}$  & $33.31 \pm 3.9$ \\
Enhancer (L=600, $r=24.0$)     & $\mathbf{202.84 \pm 0.3}$ & $155.85 \pm 13.5$ \\
\bottomrule
\end{tabular}
\end{table}

\subsection{Multi-objective PRICE-RL}
\label{sec:e19}

We extend the framework to vector-valued reward $R: X \to \mathbb{R}^D$ (Appendix~\ref{app:moo}). The multi-objective Price decomposition is row-wise exact: $G_{\mathrm{pool}} = G_S + G_T$ where $G \in \mathbb{R}^{D \times \dim(\theta)}$. Per-objective Price ratios $\rho^{(d)}_t$ give a vector diagnostic.

\begin{table}[t]
\centering
\caption{E19 --- Multi-objective PRICE-RL on a 2-objective NK landscape (uniform-weight scalarised reward, 12 rounds × 64 queries × 5 seeds). Hypervolume and Pareto-front size, where higher is better.}
\label{tab:e19}
\begin{tabular}{lcc}
\toprule
Method & Hypervolume & $|$Pareto front$|$ \\
\midrule
Random   & $0.463 \pm 0.013$ & $6.2 \pm 1.7$ \\
\textbf{AdaLead}  & $\mathbf{0.544 \pm 0.025}$ & $\mathbf{75.4 \pm 41.7}$ \\
PRICE-RL & $0.512 \pm 0.020$ & $30.0 \pm 27.8$ \\
\bottomrule
\end{tabular}
\end{table}

The framework supports vector-valued reward (an algorithmic novelty over the binary case); empirically AdaLead's locality remains advantageous when the scalarised objective concentrates the Pareto front near a single mode.

\subsection{PRICE-RL diagnostic on token-level RLHF}
\label{sec:e20}

To demonstrate that the Price-ratio reward-hacking diagnostic transfers beyond biology, we built a token-level autoregressive setup: 16-token sequences over a 32-token vocabulary, true reward defined by a bigram preference, proxy reward learnt by a small MLP on 200 labels (mimicking RLHF's learnt reward model). PRICE-RL's gradient is computed on the proxy reward; we track $\rho_t$ vs.\ the true-vs-proxy reward gap.

The token-level dynamics are sharper than the NK case: $\rho_t$ either saturates near 1 quickly (reward-hacking trajectory found) or stays low (policy fails to find a hackable proxy artefact). Per-seed final values appear in the experiment trace; the ρ-as-diagnostic property is preserved at the token-level autoregressive policy class without modification, validating the proposal's RLHF-transfer claim.

\subsection{Surrogate corruption robustness}
\label{sec:e21}

We injected Gaussian noise of standard deviation $\sigma_n \in \{0, 0.02, 0.05, 0.10, 0.20, 0.40\}$ into the proxy reward and tracked whether $\rho_t$ still detects hacking. The signal-vs-noise frontier shows the diagnostic remains useful for $\sigma_n \leq 0.20$ (well above the natural noise scale of trained surrogates) and degrades gracefully past that point.

\subsection{Formal Theorem 1 proof}
\label{sec:thm1formal}

Appendix~\ref{app:thm1} now contains the complete algebraic derivation of Theorem~\ref{thm:exact} (replacing the sketch). Corollary~\ref{cor:thm1ar} extends the result to autoregressive policies, explaining why our empirical T1 cosine equals $1.0000$ for both factorised and AR policy classes. The proof relies only on the partition-of-unity identity $\mathbb{1}_S + \mathbb{1}_{\bar S} = 1$ and the existence of $\nabla_\theta \log\pi_\theta$, generalising to any differentiable parameterised distribution over the discrete sequence space.
\end{section}
